% !TEX program = pdflatex
\documentclass[12pt,a4paper]{ctexart}

% ========== 页面设置 ==========
\usepackage[top=2.5cm, bottom=2.5cm, left=3.0cm, right=2.5cm]{geometry}
\usepackage{setspace}
\onehalfspacing
\parindent=2em

% ========== 字体 ==========
\usepackage{times}          % 英文 Times New Roman

% ========== 图表支持 ==========
\usepackage{graphicx}
\usepackage{caption}
\usepackage{float}

% ========== 超链接 ==========
\usepackage[colorlinks=true, linkcolor=black, citecolor=black, urlcolor=black]{hyperref}

% ========== 页眉页脚 ==========
\usepackage{fancyhdr}
\pagestyle{fancy}
\fancyhf{}
\fancyhead[L]{\zihao{6}\rmfamily 江苏大学课程论文}
\fancyhead[R]{\zihao{6}\rmfamily 论文标题}
\fancyfoot[C]{\zihao{6}\thepage}
\renewcommand{\headrulewidth}{0.4pt}
\renewcommand{\footrulewidth}{0pt}
\fancypagestyle{plain}{
    \fancyhf{}
    \fancyfoot[C]{\zihao{6}\thepage}
    \renewcommand{\headrulewidth}{0pt}
}
\setlength{\headheight}{14.5pt}
\addtolength{\topmargin}{-2.5pt}

% ========== 标题格式 ==========
\ctexset{
    section/format = \heiti\zihao{3},
    subsection/format = \songti\bfseries\zihao{4},
}

\begin{document}

% ==================== 封面 ====================
\thispagestyle{empty}
\begin{center}
    \vspace*{2cm}
    
    {\heiti\zihao{2} 论文题目}
    
    \vspace{3cm}
    
    {\songti\zihao{4}
    专业班级：××××\hspace{3em}学生姓名：×××
    
    \vspace{0.5cm}
    
    指导教师：×××\hspace{4.5em}职\quad 称：讲师
    }
\end{center}

\vspace{1cm}

% ==================== 中文摘要 ====================
\begin{center}
    {\heiti\zihao{4} 摘\quad 要}
\end{center}

摘要内容。摘要中不得出现引用。

\vspace{12bp}
\noindent{\heiti\zihao{4} 关键词：} 关键词1；关键词2；关键词3

% ==================== 目录 ====================
\newpage
\tableofcontents

% ==================== 正文 ====================
\newpage

\section{绪论}
正文内容，引用示例\cite{ref1}。

\subsection{研究背景与意义}
背景描述。

\section{核心技术分析}
技术分析正文。图表引用：图\ref{fig:example} 展示了……

\begin{center}
    \includegraphics[width=0.8\textwidth]{figures/fig_example.jpg}
    \captionof{figure}{图表标题}\label{fig:example}
\end{center}

\section{结论}
\subsection{主要结论}
结论内容。结论中不得出现引用。

\subsection{研究展望}
展望内容。

% ==================== 参考文献 ====================
\addcontentsline{toc}{section}{参考文献}
\renewcommand{\refname}{}
\begin{center}
    {\zihao{-2}\bfseries 参考文献}
\end{center}
\vspace{12bp}
\begin{thebibliography}{99}
\bibitem{ref1} 作者一, 作者二. 论文标题[J]. 期刊名, 年份, 卷(期): 页码.
\bibitem{ref2} 作者三. 书名[M]. 出版地: 出版社, 年份.
\end{thebibliography}

\end{document}
